%% Đề 24 -- de-thi-thu-tuyen-sinh-lop-10-toan-nam-2025-thpt-viet-bac-lang-son.pdf
\deheader{ĐỀ SỐ 24}%
  {Kỳ thi Tuyển sinh vào lớp 10 THPT — Năm 2025}%
  {Sở GD\&ĐT Lạng Sơn | Chỉ phần Đại số}
\addcontentsline{toc}{section}{Đề 24 — Tuyển sinh vào lớp 10 THPT — Năm 2025}

\begin{de_cau}{Câu 1}{2,0 điểm}
\begin{enumerate}[label=\alph*)]
  \item Tính: $A = \sqrt{12} - 2\sqrt{3} + \sqrt{27}$.
  \item Rút gọn:
  \[B = \left(\dfrac{\sqrt{x}}{\sqrt{x}+1} + \dfrac{1}{\sqrt{x}-1}\right) \cdot \dfrac{\sqrt{x}-1}{2}\]
  với $x \ge 0,\; x \neq 1$.
\end{enumerate}
\end{de_cau}

\begin{loi_giai_cau}
\textbf{a)} $A = 2\sqrt{3}-2\sqrt{3}+3\sqrt{3} = \boxed{3\sqrt{3}}$.

\textbf{b)} $\dfrac{\sqrt{x}}{\sqrt{x}+1}+\dfrac{1}{\sqrt{x}-1} = \dfrac{\sqrt{x}(\sqrt{x}-1)+(\sqrt{x}+1)}{(\sqrt{x}+1)(\sqrt{x}-1)} = \dfrac{x+1}{x-1}$.

$B = \dfrac{x+1}{x-1}\cdot\dfrac{\sqrt{x}-1}{2} = \dfrac{(x+1)(\sqrt{x}-1)}{2(x-1)} = \dfrac{x+1}{2(\sqrt{x}+1)} = \boxed{\dfrac{x+1}{2(\sqrt{x}+1)}}$.
\end{loi_giai_cau}

\begin{de_cau}{Câu 2}{2,0 điểm}
\begin{enumerate}[label=\alph*)]
  \item Giải phương trình: $x^2 - 3x - 4 = 0$.
  \item Giải hệ phương trình:
  \[\begin{cases} 2x + y = 4 \\ x - y = 2 \end{cases}.\]
\end{enumerate}
\end{de_cau}

\begin{loi_giai_cau}
\textbf{a)} $(x-4)(x+1)=0 \Rightarrow \boxed{x=4}$ hoặc $\boxed{x=-1}$.

\textbf{b)} Cộng hai pt: $3x = 6 \Rightarrow x=2, y=0$. Nghiệm $\boxed{(x;y)=(2;0)}$.
\end{loi_giai_cau}

\begin{de_cau}{Câu 3}{2,0 điểm}
Cho phương trình $x^2 - 2(m+1)x + m^2+1 = 0$ (tham số $m$).
\begin{enumerate}[label=\alph*)]
  \item Chứng minh phương trình luôn có hai nghiệm phân biệt với mọi $m$.
  \item Tìm $m$ để $x_1^2 + x_2^2 = 8$.
\end{enumerate}
\end{de_cau}

\begin{loi_giai_cau}
\textbf{a)} $\Delta' = (m+1)^2-(m^2+1)=2m>0$ khi $m>0$; và $\Delta'=(m+1)^2-(m^2+1)=2m$.

Hmm, $\Delta'=2m$, nên phương trình có 2 nghiệm phân biệt khi $m>0$.

\textbf{b)} Vi-ét: $x_1+x_2=2(m+1)$, $x_1x_2=m^2+1$.

$x_1^2+x_2^2=4(m+1)^2-2(m^2+1)=2m^2+8m+2=8 \Rightarrow m^2+4m-3=0 \Rightarrow m=-2\pm\sqrt{7}$.

Kết hợp $m>0$: $m=\boxed{-2+\sqrt{7}}$.
\end{loi_giai_cau}

\begin{de_cau}{Câu 5}{0,5 điểm}
Cho $x, y$ là các số thực dương thỏa mãn $x + y = 1$. Tìm giá trị nhỏ nhất của:
\[T = \dfrac{1}{x^2} + \dfrac{1}{y^2}.\]
\end{de_cau}

\begin{loi_giai_cau}
$T=\dfrac{1}{x^2}+\dfrac{1}{y^2} \ge \dfrac{2}{xy} \ge \dfrac{2}{\left(\frac{x+y}{2}\right)^2} = \dfrac{2}{\frac{1}{4}} = 8$.

Dấu $=$ khi $x=y=\dfrac{1}{2}$. Vậy $T_{\min} = \boxed{8}$.
\end{loi_giai_cau}
